%%%%%%%%%%  scan idx 433  |  folio 420  %%%%%%%%%%
générateurs de $J$ au voisinage de $x$ tel que le complexe de Koszul associé soit acyclique (il en existe si on suppose l'immersion régulière); soit $g_1,\ldots,g_d$ un autre système minimum de générateurs de $J$; il existe un voisinage de $x$ sur lequel $J$ peut être engendré soit par les $f_i$, soit par les $g_j$, et par suite, il existe sur ce voisinage une matrice à coefficients dans $\cO_X$ permettant d'exprimer les $g_j$ en fonction des $f_i$. Or les images des $f_i$ (resp. $g_j$) dans $J/J^2$ forment une base au point $x$, donc dans un voisinage de $x$, et par suite cette matrice est inversible dans un voisinage de $x$. Elle permet donc de définir un isomorphisme du complexe de Koszul défini par les $g_j$ sur le complexe de Koszul défini par les $f_i$, ce qui montre que le complexe de Koszul défini par les $g_j$ est acyclique.

\textbf{1.4.3.} En supposant seulement $J$ de type fini, et $J/J^2$ localement libre de type fini, notons que tout système minimum de générateurs de $J$ au voisinage d'un point $y$ de $Y$ donne une base de $J/J^2$ au voisinage de $y$. Réciproquement, si on a au voisinage de $y$ une base de $J/J^2$, on peut la relever au voisinage de $y$ en une famille de sections de $J$, qui engendrent $J_y$ par Nakayama, donc $J$ dans un voisinage de $y$ puisque $J$ est de type fini; de plus ces générateurs sont en nombre minimum, puisqu'ils donnent une base de $J/J^2$.

\textbf{Proposition 1.5.} (Comparer EGA IV 19.1.5 (ii)). Soit $i:Y\longrightarrow X$ une immersion. Si $i$ est régulière, il en est de même pour toute immersion déduite de $i$ par changement de base plat. Pour que $i$ soit régulière, il faut et il suffit qu'elle le soit localement pour la topologie fidèlement plate quasi-compacte.

La stabilité par changement de base plat est évidente.

Réciproquement, soit $i'$ l'immersion déduite de $i$ par un changement de
